\begin{abstract}
Rolling shutter (RS) cameras are widely deployed in robotic systems due to their low cost and power efficiency. However, robot ego-motion induces significant RS distortions that degrade geometric consistency and severely challenge visual perception. In this work, we address monocular depth estimation under large rolling shutter distortions caused by ego-motion. We introduce the concept of an anchor point, which establishes a consistent temporal and geometric reference between rolling shutter and global shutter observations. This formulation enables principled adaptation and fine-tuning of deep foundation models for depth estimation under RS distortions. To systematically evaluate the problem, we present a new rolling shutter dataset featuring large motion-induced geometric distortions representative of mobile robotic tasks. Our experimental analysis shows that anchor point selection is a primary factor governing depth accuracy under RS effects. We further characterize anchor configurations for RS-GS data capture and identify stable regimes that yield accurate depth estimates. Using a representative two-stage depth estimation pipeline, we conduct extensive ablation studies demonstrating that anchor-based modeling significantly improves performance, outperforming recent RS correction baselines under strong ego-motion. Finally, we validate our approach on a legged robot platform, demonstrating robust monocular depth estimation in the presence of substantial rolling shutter distortions in under 400 ms.
%In this paper, we focus on monocular depth estimation, a crucial robot perception task, in the face of large RS distortions. We introduce the concept of an \textit{anchor point}, which defines a consistent temporal and geometric reference between rolling shutter and global shutter observations and enables effective fine-tuning of deep foundational models under RS distortions. We also introduce a new RS dataset to evaluate our method with large geometric distortions that are common in mobile robot vision tasks. Our experiments show that the choice of temporal anchor point enables high performance on our new dataset, where egocentric motions amplify RS distortions. We analyze anchor point selection for RS-GS data capture and identify configurations that yield stable and accurate depth estimates. Using a representative two-stage depth estimation pipeline, we show through extensive ablations that anchor point choice and dataset design are primary drivers of performance, outperforming recent RS correction baselines under large ego-motion. We also show a legged robot experiment demonstrating near real-time performance of the monocular depth estimation.
\end{abstract}